\chapter{总结与展望}
\label{chap:conclusion}

\section{全文总结}

本文从「同样的文字，排法不同，读起来累不累也不同」这一朴素观察出发，
提出了版面负荷指标，并把它与语言特征融合为一个可读性评估模型。
实验表明，版面因素可以解释可读性差异中约五分之一的部分，
在信息密集的技术文档中尤为明显。

\section{不足与展望}

\begin{enumerate}
  \item 实验只覆盖了横排简体中文，竖排与繁体文本有待检验；
  \item 版面负荷指标没有考虑字体本身的差异，后续可以引入字重、字面大小等特征；
  \item 疲劳程度依赖读者自评，将来可以结合眼动数据做更细致的测量。
\end{enumerate}
